\chapter{Conclusions and Future Work}
 
\section{Conclusions}

This thesis studied the use of Bayesian Adaptive Spline Surfaces (BASS) as surrogate models within a Bayesian Optimization framework for black-box optimization and hyperparameter tuning. The central objective was to determine whether a BASS-based optimizer can offer competitive sample efficiency when compared with established baselines, while also providing practical flexibility for difficult objective landscapes. \green{To evaluate this, a complete experimental pipeline was developed in R, including multi-dimensional optimization support, reproducible seed-based replication, parallel execution, and standardized convergence and paired-statistics reporting against GP-BO, Random Search, and TPE.}

% PASS2: replace the conditional phrasing below with the measured outcomes from
% final_results/ (continuous benchmarks: relative convergence of GP-BO and
% BASS-BO; note whichever way it comes out).
\green{The experimental chapter reports, for every benchmark, mean convergence curves with confidence intervals and per-seed paired comparisons against the Random Search baseline. On the smooth, low-dimensional continuous benchmarks, where GP surrogates are particularly well matched to the objective, GP-BO is expected to converge fastest; the results chapter reports whether that expectation holds and by how much. Random Search provides the floor that any surrogate-guided method must clear, and the paired win/tie/loss counts make that comparison explicit seed by seed.}

\textbf{A key conclusion of this work is that BASS-BO should not be interpreted as a universally superior replacement for GP-BO.} Instead, the two should be seen as suited to different regimes. \green{On the smooth, low-dimensional continuous benchmarks, GP-BO is the natural favorite, both because its smoothness assumptions match these problems and because its closed-form predictive uncertainty is ideally matched to them; the experimental chapter quantifies how large that advantage actually is.} Importantly, in this implementation the two methods are placed on equal footing: they share the same parameter-free Expected Improvement acquisition and the same candidate generator, and BASS supplies its predictive uncertainty natively through its posterior draws rather than through any hand-tuned exploration schedule. The difference in performance is therefore attributable to the surrogate itself rather than to tuning, which is exactly what makes the comparison clean.

From a methodological perspective, one contribution of this thesis is the construction of a reproducible comparison in which the surrogate is the only moving part. Holding the full experimental protocol constant, including the initialization strategy, candidate pool size, duplicate handling, seed replication, and stopping budget, is what allows measured differences to be read as differences in model quality rather than in procedure. Without this consistency, apparent performance gaps can reflect the surrounding machinery rather than the surrogate.

\green{The development of this pipeline produced direct evidence for how large that machinery effect can be, and we consider this a finding in its own right. In controlled oracle experiments --- replacing the acquisition with the true objective, so that the candidate pool is the only binding constraint --- a single detail of the candidate generator (whether local proposals may keep the incumbent's categorical combination while refining its continuous coordinates) determined whether \emph{any} surrogate, including a perfect one, could approach the optimum of the mixed benchmarks: the restricted generator was beaten by the permissive one in every paired seed, and with the restricted generator even a perfect surrogate exploited no faster than the pool allowed. Likewise, before duplicate handling operated at the level of decoded categorical combinations, roughly two thirds of the evaluation budget on a small categorical benchmark was silently spent re-evaluating known combinations. Both effects are invisible in a standard convergence plot, and both would have been attributed to the surrogate had the machinery not been shared. The implication for the field is uncomfortable but useful: surrogate comparisons in mixed and categorical spaces are confounded by acquisition-optimization machinery unless that machinery is shared, categorical-aware, and audited --- and much of the published literature compares surrogates atop method-specific machinery. A systematic treatment of this confound, including the oracle-ceiling audit as a general diagnostic, is beyond the scope of this thesis and is being developed as a separate article.}

\green{The current study also has scope limits. The benchmark set spans two to six input dimensions and a single evaluation budget, and the continuous benchmarks in particular emphasize smooth functions where GP assumptions are favorable. The categorical and mixed-variable benchmarks probe the opposite regime; the Cat-Ackley family is deliberately run at three sizes so that capability (on a space the budget can cover) and scaling behavior (on spaces that outgrow it) are reported separately rather than conflated.} \textbf{Broader coverage, including more irregular, nonstationary, or interaction-heavy objectives}, remains an open empirical direction that the established codebase is now prepared to test.
% PASS2: if the final numbers are in, state the categorical outcomes here
% (easy-instance solve rates, hard-instance paired wins vs Random, and the
% Func-2C/3C result -- including parity with Random if that is what held).

The practical takeaway is therefore one of complementary strengths. \green{For smooth, low-dimensional objectives under tight evaluation budgets, GP-BO is an excellent default; the experimental chapter reports how the two compare there under the shared protocol.}
% PASS2: state the measured continuous-benchmark outcome in the sentence above.
\textbf{The distinguishing advantage of BASS-BO is structural rather than a matter of tuning: it accepts categorical variables natively}, through basis functions defined on subsets of their levels, and can therefore operate on mixed and purely categorical search spaces on which a standard Gaussian process must fall back on encodings or specialized kernels. This is the capability that the categorical benchmarks of the previous chapter are designed to exercise, and it is where BASS-BO is positioned to be the more natural choice. It is worth being clear that BASS is not the only surrogate with this property: tree- and density-based methods, the Tree-structured Parzen Estimator chief among them, also handle categorical inputs natively, which is precisely why we benchmark against it rather than against a Gaussian process alone. The contribution is therefore not that BASS is uniquely categorical-capable, but that it brings a flexible, fully Bayesian spline surrogate to a setting the Gaussian process it is usually compared against cannot reach, and does so competitively with a purpose-built categorical optimizer.

\section{Future Work}

\green{Future work should extend this analysis in two directions. First, the categorical and mixed-variable evaluation should be deepened: scaling the number of categories and categorical dimensions further, enriching the categorical proposal distribution beyond the current per-coordinate Hamming moves (for example, correlated multi-coordinate moves informed by the surrogate's own level effects), and moving from synthetic benchmarks to real mixed hyperparameter-tuning tasks, where categorical choices such as model family or kernel sit alongside continuous parameters.} Second, the continuous benchmark suite should be broadened to include more multimodal, higher-dimensional, and noisy objectives, so that conclusions are less tied to any one function class. Together these would sharpen the characterization of when BASS-BO is preferable to GP-BO.

% PASS2: finalize this summary with the measured outcomes (continuous and
% categorical) once final_results/ is committed.
\green{In summary, this thesis delivers a reproducible and extensible framework for comparing surrogate-based Bayesian Optimization methods, and demonstrates that BASS can be integrated effectively into a multi-dimensional BO loop. The central contribution is to identify and operationalize the regime where BASS-BO is genuinely distinct: optimization over categorical and mixed search spaces, which a standard Gaussian process cannot address natively. The main outcome is therefore both empirical and methodological: the experimental chapter quantifies each method's standing on the smooth continuous problems, while BASS-BO extends surrogate-based Bayesian Optimization to the categorical setting within a single, consistent framework.}
